% French Study Book - Main Document
% French Study Book - LaTeX Preamble
% Encoding and language settings
\documentclass[11pt, a4paper, twoside, openright]{book}
\usepackage[utf8]{inputenc}
\usepackage[T1]{fontenc}
\usepackage[english, french]{babel}
\usepackage[autoref, colorlinks=true, linkcolor=blue, citecolor=blue, urlcolor=blue, pdfborder={0 0 0}]{hyperref}

% Fonts - Using system fonts for better embedding
\usepackage{fontspec}
\setmainfont{Linux Libertine O}
\newfontfamily\headingfont{Playfair Display}
\newfontfamily\sansfont{Source Sans Pro}

% Page geometry
\usepackage[a4paper, inner=25mm, outer=20mm, top=25mm, bottom=25mm, bindingoffset=5mm]{geometry}

% Colors - French-inspired palette
\definecolor{frenchblue}{RGB}{0, 85, 164}
\definecolor{frenchred}{RGB}{232, 17, 45}
\definecolor{frenchgold}{RGB}{255, 215, 0}
\definecolor{cream}{RGB}{245, 245, 220}
\definecolor{darkgray}{RGB}{44, 44, 44}
\definecolor{lightgray}{RGB}{200, 200, 200}

% Typographic enhancements
\usepackage{microtype}
\usepackage{setspace}
\setstretch{1.5}
\usepackage{parskip}
\setlength{\parskip}{1em}
\setlength{\parindent}{0pt}

% Chapter and section styling
\usepackage{titlesec}
\usepackage{titletoc}

% Custom chapter style
\titleformat{\chapter}[display]
    {\normalfont\headingfont\huge\bfseries\color{frenchblue}}
    {\chaptertitlename\ \thechapter}
    {20pt}
    {\Huge\color{darkgray}}
    [\vspace{-10pt}{\color{frenchred}\rule{\textwidth}{2pt}}]

\titleformat{name=\chapter,numberless}[display]
    {\normalfont\headingfont\huge\bfseries\color{frenchblue}}
    {}
    {0pt}
    {\Huge\color{darkgray}}
    [\vspace{-10pt}{\color{frenchred}\rule{\textwidth}{2pt}}]

% Section styling
\titleformat{\section}
    {\normalfont\headingfont\Large\bfseries\color{frenchblue}}
    {\thesection}
    {1em}
    {\color{darkgray}}
    [\vspace{-5pt}{\color{frenchred}\rule{0.5\textwidth}{1pt}}]

\titleformat{\subsection}
    {\normalfont\headingfont\large\bfseries\color{frenchblue}}
    {\thesubsection}
    {1em}
    {\color{darkgray}}

% Page headers and footers
\usepackage{fancyhdr}
\pagestyle{fancy}
\fancyhf{}
\fancyhead[LE,RO]{\thepage}
\fancyhead[RE]{\nounchap\leftmark}
\fancyhead[LO]{\rightmark}
\renewcommand{\headrulewidth}{2pt}
\renewcommand{\headrule}{\hbox to\headwidth{\color{frenchred}\leaders\hrule height \headrulewidth\hfill}}
\renewcommand{\footrulewidth}{0pt}

% For plain pages (chapter starts)
\fancypagestyle{plain}{
    \fancyhf{}
    \fancyfoot[C]{\thepage}
    \renewcommand{\headrulewidth}{0pt}
}

% Tables
\usepackage{tabularx}
\usepackage{booktabs}
\usepackage{multirow}
\usepackage{longtable}
\usepackage{array}
\newcolumntype{L}{>{\raggedright\arraybackslash}X}
\newcolumntype{R}{>{\raggedleft\arraybackslash}X}
\newcolumntype{C}{>{\centering\arraybackslash}X}

% Lists
\usepackage{enumitem}
\setlist[itemize]{leftmargin=*, topsep=0.5em, itemsep=0.25em}
\setlist[enumerate]{leftmargin=*, topsep=0.5em, itemsep=0.25em}

% Boxes for special content
\usepackage{tcolorbox}
\tcbuselibrary{most}
\newtcolorbox{mascotbox}[1][]{
    colback=cream,
    colframe=frenchblue,
    arc=5pt,
    boxrule=2pt,
    left=10pt,
    right=10pt,
    top=8pt,
    bottom=8pt,
    fonttitle=\headingfont\bfseries\large,
    title=#1,
    coltitle=frenchblue,
    enhanced,
    breakable
}

\newtcolorbox{dialoguebox}{
    colback=white,
    colframe=frenchred,
    arc=5pt,
    boxrule=1.5pt,
    left=10pt,
    right=10pt,
    top=8pt,
    bottom=8pt,
    enhanced,
    breakable
}

\newtcolorbox{vocabbox}{
    colback=white,
    colframe=frenchgold,
    arc=5pt,
    boxrule=1.5pt,
    left=10pt,
    right=10pt,
    top=8pt,
    bottom=8pt,
    enhanced,
    breakable
}

\newtcolorbox{exercisebox}[1][]{
    colback=white,
    colframe=frenchblue,
    arc=5pt,
    boxrule=1.5pt,
    left=10pt,
    right=10pt,
    top=8pt,
    bottom=8pt,
    fonttitle=\headingfont\bfseries,
    title=#1,
    coltitle=frenchblue,
    enhanced,
    breakable
}

\newtcolorbox{notebox}{
    colback=cream,
    colframe=frenchred,
    arc=5pt,
    boxrule=1.5pt,
    left=10pt,
    right=10pt,
    top=8pt,
    bottom=8pt,
    fonttitle=\headingfont\bfseries,
    title=Note,
    coltitle=frenchred,
    enhanced,
    breakable
}

\newtcolorbox{culturebox}{
    colback=white,
    colframe=frenchgold,
    arc=5pt,
    boxrule=1.5pt,
    left=10pt,
    right=10pt,
    top=8pt,
    bottom=8pt,
    fonttitle=\headingfont\bfseries,
    title=Cultural Note,
    coltitle=frenchgold,
    enhanced,
    breakable
}

% SVG graphics
\usepackage{svg}
\svgpath{{/workspace/French-Study-Book/content/assets/svg/}}

% Floating figures
\usepackage{float}
\floatstyle{plaintop}
\newfloat{illustration}{htbp}{lop}
\floatname{illustration}{Illustration}

% Custom environments
\newenvironment{mascot}[1][]{
    \begin{center}
    \begin{mascotbox}[#1]
    \includesvg[height=8cm]{#1}
}{
    \end{mascotbox}
    \end{center}
}

% French-specific typography
\usepackage{frenchle}
\frenchbsetup{
    FrenchFootnotes=false,
    og=«",
    fg=»",
    thincolonspace=false
}

% Mathematics (for any linguistic diagrams)
\usepackage{amsmath}
\usepackage{amssymb}

% Index
\usepackage{makeidx}
\makeindex

% Glossary
\usepackage{glossaries}
\makeglossaries

% Custom commands
\newcommand{\frenchword}[1]{\textbf{#1}}
\newcommand{\pronunciation}[1]{/\textit{#1}/}
\newcommand{\grammar}[1]{\textcolor{frenchblue}{\textbf{#1}}}
\newcommand{\vocabulary}[1]{\textcolor{frenchgold}{\textbf{#1}}}
\newcommand{\dialogue}[2]{\begin{dialoguebox}\textbf{#1:} #2\end{dialoguebox}}
\newcommand{\level}[1]{\textcolor{white}\fcolorbox{frenchblue}{\parbox{\linewidth}{\centering\Large\bfseries CEFR Level #1}}}

% Level indicators
\newcommand{\levelAone}{\textcolor{white}\fcolorbox{frenchblue}{\parbox{1cm}{\centering\bfseries A1}}}
\newcommand{\levelAtwo}{\textcolor{white}\fcolorbox{frenchred}{\parbox{1cm}{\centering\bfseries A2}}}
\newcommand{\levelBone}{\textcolor{white}\fcolorbox{frenchgold}{\parbox{1cm}{\centering\bfseries B1}}}
\newcommand{\levelBtwo}{\textcolor{white}\fcolorbox{frenchblue}{\parbox{1cm}{\centering\bfseries B2}}}
\newcommand{\levelCone}{\textcolor{white}\fcolorbox{frenchred}{\parbox{1cm}{\centering\bfseries C1}}}
\newcommand{\levelCtwo}{\textcolor{white}\fcolorbox{frenchgold}{\parbox{1cm}{\centering\bfseries C2}}}

% Answer key formatting
\newcommand{\answer}[1]{\textcolor{frenchblue}{\textbf{Answer:} #1}}

% Exercise types
\newcommand{\listening}{\textcolor{frenchred}{\textbf{Listening}}}
\newcommand{\speaking}{\textcolor{frenchgold}{\textbf{Speaking}}}
\newcommand{\reading}{\textcolor{frenchblue}{\textbf{Reading}}}
\newcommand{\writing}{\textcolor{frenchred}{\textbf{Writing}}}
\newcommand{\grammarex}{\textcolor{frenchgold}{\textbf{Grammar}}}
\newcommand{\vocabex}{\textcolor{frenchblue}{\textbf{Vocabulary}}}

% PDF metadata
\hypersetup{
    pdftitle={Complete French: A Self-Study Textbook},
    pdfauthor={Amélie - Your French Learning Companion},
    pdfsubject={French Language Learning - CEFR A1 to C2},
    pdfkeywords={French, language, learning, self-study, CEFR, A1, A2, B1, B2, C1, C2},
    pdfcreator={XeLaTeX},
    pdfproducer={French Study Book Project},
    colorlinks=true,
    linkcolor=frenchblue,
    citecolor=frenchred,
    urlcolor=frenchgold,
    bookmarksopen=true,
    bookmarksnumbered=true,
    bookmarksopenlevel=2
}

% Custom title page
\newcommand{\makecover}{
    \thispagestyle{empty}
    \begin{center}
    \vspace*{2cm}
    {\headingfont\Huge\bfseries\color{frenchblue} Complete French\par}
    \vspace{0.5cm}
    {\headingfont\huge\bfseries A Self-Study Textbook\par}
    \vspace{1cm}
    {\headingfont\LARGE\color{frenchred} From Beginner to Proficiency\par}
    \vspace{0.5cm}
    {\headingfont\Large\color{darkgray} CEFR Levels A1 to C2\par}
    \vspace{2cm}
    \includesvg[height=10cm]{mascot_amelie_welcome}
    \vspace{1cm}
    {\headingfont\Large\bfseries With Amélie, Your French Learning Companion\par}
    \vfill
    {\headingfont\large\color{darkgray} First Edition, 2024\par}
    
    \clearpage
    \thispagestyle{empty}
    \vspace*{2cm}
    {\headingfont\Large\bfseries Complete French\par}
    {\headingfont\large A Self-Study Textbook\par}
    \vspace{1cm}
    {\headingfont\normalsize\color{darkgray} CEFR Levels A1 to C2\par}
    \vspace{2cm}
    \begin{flushright}
    {\headingfont\normalsize\bfseries Written and Designed by:\par}
    {\headingfont\normalsize Amélie - Your French Learning Companion\par}
    
    \vspace{1cm}
    {\headingfont\small\color{darkgray} This book is designed for independent learners who want to\par}
    {\headingfont\small\color{darkgray} achieve fluency in French from complete beginner to\par}
    {\headingfont\small\color{darkgray} advanced proficiency.\par}
    \end{flushright}
    
    \clearpage
}

% Level introduction page
\newcommand{\levelintro}[3]{
    \clearpage
    \thispagestyle{empty}
    \begin{center}
    \vspace*{2cm}
    {\headingfont\Huge\bfseries\color{frenchblue} #1\par}
    \vspace{0.5cm}
    {\headingfont\Large\color{frenchred} #2\par}
    \vspace{1cm}
    \includesvg[height=8cm]{#3}
    \vspace{1cm}
    \begin{minipage}{0.8\textwidth}
    \centering
    {\headingfont\normalsize #2\par}
    
    \vspace{0.5cm}
    \end{minipage}
    \vfill
    
    \clearpage
}

% Unit header
\newcommand{\unitheader}[3]{
    \clearpage
    \begin{center}
    \vspace*{1cm}
    {\headingfont\huge\bfseries\color{frenchblue} Unit #1\par}
    \vspace{0.5cm}
    {\headingfont\Large\color{darkgray} #2\par}
    \vspace{0.5cm}
    {\headingfont\normalsize\color{frenchred} #3\par}
    \vspace{0.5cm}
    \includesvg[height=5cm]{mascot_amelie_explaining}
    \vspace{1cm}
    \end{center}
}

% Learning objectives
\newcommand{\objectives}[1]{
    \begin{tcolorbox}[colback=cream, colframe=frenchblue, arc=5pt, boxrule=2pt, title=\headingfont\bfseries\large Learning Objectives]
    \begin{itemize}
    #1
    \end{itemize}
    \end{tcolorbox}
}

% Vocabulary list
\newcommand{\vocabularylist}[1]{
    \begin{vocabbox}
    \begin{longtable}{l l l}
    \textbf{French} & \textbf{English} & \textbf{Pronunciation} \\
    \midrule
    #1
    \end{longtable}
    \end{vocabbox}
}

% Dialogue
\newcommand{\conversation}[1]{
    \begin{dialoguebox}
    \begin{itemize}[leftmargin=*, topsep=0.25em, itemsep=0.1em]
    #1
    \end{itemize}
    \end{dialoguebox}
}

% Exercise
\newcommand{\exercise}[2][]{
    \begin{exercisebox}[#1]
    #2
    \end{exercisebox}
}

% Answer key
\newcommand{\answerkey}[1]{
    \begin{tcolorbox}[colback=white, colframe=frenchred, arc=5pt, boxrule=1.5pt, title=\headingfont\bfseries Answer Key]
    #1
    \end{tcolorbox}
}

% Grammar explanation
\newcommand{\grammarexplanation}[2][]{
    \begin{tcolorbox}[colback=white, colframe=frenchgold, arc=5pt, boxrule=1.5pt, title=\headingfont\bfseries #1]
    #2
    \end{tcolorbox}
}

% Cultural note
\newcommand{\culturalnote}[1]{
    \begin{culturebox}
    #1
    \end{culturebox}
}

% Revision section
\newcommand{\revision}[1]{
    \begin{tcolorbox}[colback=cream, colframe=frenchblue, arc=5pt, boxrule=2pt, title=\headingfont\bfseries\large Revision]
    #1
    \end{tcolorbox}
}

% Assessment
\newcommand{\assessment}[1]{
    \begin{tcolorbox}[colback=white, colframe=frenchred, arc=5pt, boxrule=2pt, title=\headingfont\bfseries\large Assessment]
    #1
    \end{tcolorbox}
}


\begin{document}

% Cover and title pages
\makecover

% Copyright page
\thispagestyle{empty}
\begin{center}
\vspace*{2cm}
{\headingfont\Large\bfseries Complete French: A Self-Study Textbook\par}
\vspace{1cm}
{\headingfont\normalsize Copyright \textcopyright\ 2024 Amélie - French Learning Companion\par}
\vspace{0.5cm}
{\headingfont\small All rights reserved.\par}
\vspace{0.5cm}
{\headingfont\small No part of this publication may be reproduced, distributed, or transmitted\par}
{\headingfont\small in any form or by any means without prior written permission.\par}
\vfill
{\headingfont\small First Edition, October 2024\par}
{\headingfont\small Printed in digital format\par}
\end{center}

\clearpage

% Table of Contents
\pagenumbering{roman}
\setcounter{page}{1}
\tableofcontents
\markboth{Table of Contents}{Table of Contents}

\clearpage

% List of Illustrations
\listoffigures
\markboth{List of Illustrations}{List of Illustrations}

\clearpage

% List of Tables
\listoftables
\markboth{List of Tables}{List of Tables}

\clearpage

% Preface
\chapter*{Preface}
\markboth{Preface}{Preface}
\addcontentsline{toc}{chapter}{Preface}

\begin{center}
\includesvg[height=6cm]{mascot_amelie_welcome}
\end{center}

Bonjour, dear learner! I'm Amélie, your French learning companion, and I'm so excited to embark on this journey with you.

This book is designed to take you from complete beginner to advanced proficiency in French, following the Common European Framework of Reference for Languages (CEFR). Whether you're learning for travel, work, study, or simply for the love of the language, this comprehensive guide will support you every step of the way.

\begin{itemize}
\item \textbf{Complete Coverage:} All six CEFR levels (A1 to C2) with substantial content for each
\item \textbf{Independent Learning:} Everything you need for self-study, including exercises, answer keys, and assessments
\item \textbf{Integrated Culture:} French culture and history woven naturally into language learning
\item \textbf{Clear Progression:} Carefully structured units that build on each other
\item \textbf{Professional Design:} Beautiful, readable layout with illustrations throughout
\end{itemize}

Remember, learning a language is a journey, not a race. Take your time, enjoy the process, and don't be afraid to make mistakes—they're a natural part of learning!

\begin{flushright}
\textit{Bonne chance!} \par
Amélie
\end{flushright}

\clearpage

% How to Use This Book
\chapter*{How to Use This Book}
\markboth{How to Use This Book}{How to Use This Book}
\addcontentsline{toc}{chapter}{How to Use This Book}

This book is organized into six parts, corresponding to the CEFR levels: A1, A2, B1, B2, C1, and C2. Each level contains multiple units with the following structure:

\begin{itemize}
\item \textbf{Contextual Dialogue or Reading:} Real-life conversations or texts to introduce new vocabulary and grammar in context
\item \textbf{Vocabulary in Context:} New words and phrases with English translations and pronunciation guides
\item \textbf{Accessible Explanations:} Clear grammar explanations with examples
\item \textbf{Varied Exercises:} Practice activities for listening, speaking, reading, and writing
\item \textbf{Practical Communication Tasks:} Real-world applications of what you've learned
\item \textbf{Cultural Notes:} Insights into French culture, history, and society
\end{itemize}

\subsection*{Study Tips}
\begin{itemize}
\item \textbf{Set Regular Goals:} Aim for consistent, manageable study sessions
\item \textbf{Practice All Skills:} Don't neglect speaking and listening
\item \textbf{Review Regularly:} Use the revision sections to reinforce learning
\item \textbf{Immerse Yourself:} Surround yourself with French through media, music, and conversation
\item \textbf{Be Patient:} Language learning takes time and practice
\end{itemize}

\subsection*{Suggested Learning Schedules}

\begin{table}[h]
\centering
\begin{tabular}{l l l l}
\toprule
\textbf{Intensity} & \textbf{Time Commitment} & \textbf{Duration} & \textbf{Expected Progress} \\
\midrule
Intensive & 3-4 hours/day & 6-8 months & A1 to B1 \\
Regular & 1-2 hours/day & 12-18 months & A1 to B2 \\
Casual & 30-60 minutes/day & 24-36 months & A1 to C1 \\
\bottomrule
\end{tabular}
\caption{Learning Schedule Options}
\label{tab:schedule}
\end{table}

\clearpage

% Mascot Introduction
\chapter*{Meet Amélie, Your Learning Companion}
\markboth{Meet Amélie}{Meet Amélie}
\addcontentsline{toc}{chapter}{Meet Amélie, Your Learning Companion}

\begin{center}
\includesvg[height=8cm]{mascot_amelie_welcome}
\end{center}

\begin{mascotbox}[Amélie's Story]
\begin{itemize}
\item \textbf{Name:} Amélie Dubois
\item \textbf{Age:} 34
\item \textbf{From:} Lyon, France
\item \textbf{Occupation:} French Language Teacher and Cultural Guide
\item \textbf{Personality:} Friendly, patient, encouraging, and passionate about sharing French language and culture
\item \textbf{Interests:} Literature, art, cooking, travel, and helping others learn
\end{itemize}
\end{mascotbox}

Throughout this book, you'll see me in different poses:

\begin{itemize}
\item \textbf{Welcome Pose:} Introducing new topics and units
\item \textbf{Explaining Pose:} Clarifying grammar points and tricky concepts
\item \textbf{Encouraging Pose:} Celebrating your progress and offering motivation
\end{itemize}

I'll be with you every step of the way, offering explanations, encouragement, and cultural insights. Think of me as your personal French tutor, always ready to help!

\clearpage

% ============================================================================
% LEVEL A1 - BEGINNER
% ============================================================================

\levelintro{A1}{Beginner}{mascot_amelie_welcome}

\pagenumbering{arabic}
\setcounter{page}{1}

\chapter{A1: Getting Started with French}
\label{chap:a1}

\section*{About Level A1}
\markboth{Level A1: Getting Started}{Level A1: Getting Started}

\levelAone\ At the end of Level A1, you will be able to:

\begin{itemize}
\item Understand and use familiar everyday expressions and very basic phrases
\item Introduce yourself and others and ask and answer questions about personal details
\item Interact in a simple way provided the other person talks slowly and clearly
\end{itemize}

\begin{notebox}
Level A1 is for complete beginners. Don't worry if everything seems new at first—it will become familiar with practice!
\end{notebox}

\clearpage

% A1 Unit 01: Introduction and Greetings
\unitheader{1}{Introduction and Greetings}{Meeting people and basic greetings}

\objectives{
\item Learn to greet people in French
\item Introduce yourself and others
\item Understand and use basic courtesy phrases
\item Recognize French alphabet pronunciation
}

\section{Dialogue: Meeting Someone New}

\conversation{
\item \textbf{Marie:} Bonjour! Comment vous appelez? \par \textit{Hello! What's your name?}
\item \textbf{Paul:} Bonjour! Je m'appelle Paul. Et vous? \par \textit{Hello! My name is Paul. And you?}
\item \textbf{Marie:} Je m'appelle Marie. Enchantée! \par \textit{My name is Marie. Nice to meet you!}
\item \textbf{Paul:} Enchanté! Comment ça va? \par \textit{Nice to meet you! How are you?}
\item \textbf{Marie:} Ça va bien, merci. Et vous? \par \textit{I'm fine, thank you. And you?}
\item \textbf{Paul:} Très bien, merci. \par \textit{Very well, thank you.}
\item \textbf{Marie:} Au revoir! \par \textit{Goodbye!}
\item \textbf{Paul:} Au revoir! \par \textit{Goodbye!}
}

\section{Vocabulary}

\vocabularylist{
01_introduction & Bonjour & Hello & \pronunciation{bong-zhoor} \\
 & Bonsoir & Good evening & \pronunciation{bong-swahr} \\
 & Salut & Hi (informal) & \pronunciation{sa-loo} \\
 & Au revoir & Goodbye & \pronunciation{oh rev-wahr} \\
 & À bientôt & See you soon & \pronunciation{a bee-en-toh} \\
 & Comment vous appelez? & What's your name? (formal) & \pronunciation{kom-mong voo za-play} \\
 & Je m'appelle... & My name is... & \pronunciation{zhe ma-pehl} \\
 & Et vous? & And you? & \pronunciation{ay voo} \\
 & Enchanté(e) & Nice to meet you & \pronunciation{on-shon-tay} \\
 & Comment ça va? & How are you? & \pronunciation{kom-mong sa va} \\
 & Ça va bien & I'm fine & \pronunciation{sa va bee-en} \\
 & Très bien & Very well & \pronunciation{tray bee-en} \\
 & Merci & Thank you & \pronunciation{mehr-see} \\
}

\section{Grammar: Basic Sentence Structure}

\grammarexplanation[Subject Pronouns and "to be"]{
In French, the verb \frenchword{être} (to be) is essential. Here are the subject pronouns and the present tense of \frenchword{être}:

\begin{table}[h]
\centering
\begin{tabular}{l l l}
\toprule
\textbf{French} & \textbf{English} & \textbf{Pronunciation} \\
\midrule
je suis & I am & \pronunciation{zhe swee} \\
 tu es & you are (singular, informal) & \pronunciation{too ay} \\
 il/elle/on est & he/she/one is & \pronunciation{eel/ell/on ay} \\
 nous sommes & we are & \pronunciation{noo som} \\
 vous êtes & you are (plural/formal) & \pronunciation{voo zet} \\
 ils/elles sont & they are & \pronunciation{eel/ell sont} \\
\bottomrule
\end{tabular}
\caption{Present Tense of \frenchword{Être} (to be)}
\label{tab:etre}
\end{table}

\begin{itemize}
\item \frenchword{Je m'appelle Paul} = I call myself Paul = My name is Paul
\item \frenchword{Comment vous appelez?} = How do you call yourself? = What's your name?
\end{itemize}
}

\section{Exercises}

\exercise[Vocabulary Practice]{
Match the French words with their English meanings:

\begin{table}[h]
\centering
\begin{tabular}{l l}
\toprule
\textbf{French} & \textbf{English} \\
\midrule
1. Bonjour & a. Goodbye \\
2. Au revoir & b. Hello \\
3. Merci & c. Thank you \\
4. Comment ça va? & d. How are you? \\
\bottomrule
\end{tabular}
\end{table}
}

\answerkey{
1. b, 2. a, 3. c, 4. d
}

\exercise[Listening Practice]{
\listening\ Read the following dialogue aloud, then listen to it using a text-to-speech tool:

\begin{itemize}
\item Marie: Bonjour! Comment vous appelez?
\item Paul: Bonjour! Je m'appelle Paul. Et vous?
\item Marie: Je m'appelle Marie. Enchantée!
\item Paul: Enchanté! Comment ça va?
\item Marie: Ça va bien, merci. Et vous?
\item Paul: Très bien, merci.
\end{itemize}

Then answer: What is Paul's name? What is Marie's name?
}

\answerkey{
Paul's name is Paul. Marie's name is Marie.
}

\exercise[Speaking Practice]{
\speaking\ Practice the dialogue with a partner or record yourself. Try to sound natural and confident!
}

\exercise[Writing Practice]{
\writing\ Write a short dialogue (4-6 lines) introducing yourself to someone new. Use the phrases you've learned.
}

\section{Cultural Note}

\culturalnote{
In France, it's customary to greet people with \frenchword{Bonjour} (during the day) or \frenchword{Bonsoir} (in the evening). When entering a shop, restaurant, or any business, always greet the staff first. This is considered polite and is expected in French culture.

Also, the French often use \frenchword{Comment ça va?} (How are you?) as a casual greeting, similar to "How's it going?" in English. The typical response is \frenchword{Ça va bien, merci} (I'm fine, thank you).
}

\section{Revision}

\revision{
\begin{itemize}
\item Practice greeting someone and introducing yourself
\item Memorize the subject pronouns and the verb \frenchword{être}
\item Review the vocabulary list
\item Try the dialogue with different names
\end{itemize}
}

\clearpage

% A1 Unit 02: Numbers and Alphabet
\unitheader{2}{Numbers and Alphabet}{Counting and spelling in French}

\objectives{
\item Learn to count from 0 to 100
\item Understand and use the French alphabet
\item Spell your name and other words
\item Ask and answer questions about quantities
}

\section{Dialogue: At the Market}

\conversation{
\item \textbf{Vendeur:} Bonjour! Je peux vous aider? \par \textit{Hello! Can I help you?}
\item \textbf{Client:} Bonjour! Je voudrais trois pommes, s'il vous plaît. \par \textit{Hello! I would like three apples, please.}
\item \textbf{Vendeur:} Trois pommes? Et avec ça? \par \textit{Three apples? And with that?}
\item \textbf{Client:} Deux bananes aussi, s'il vous plaît. \par \textit{Two bananas as well, please.}
\item \textbf{Vendeur:} Ça fait cinq euros. \par \textit{That makes five euros.}
\item \textbf{Client:} Voici dix euros. \par \textit{Here is ten euros.}
\item \textbf{Vendeur:} Merci! Au revoir! \par \textit{Thank you! Goodbye!}
\item \textbf{Client:} Au revoir! \par \textit{Goodbye!}
}

\section{Vocabulary}

\vocabularylist{
02_numbers_alphabet & zéro & zero & \pronunciation{zay-ro} \\
 & un & one & \pronunciation{uhn} \\
 & deux & two & \pronunciation{duh} \\
 & trois & three & \pronunciation{trwah} \\
 & quatre & four & \pronunciation{ka-truh} \\
 & cinq & five & \pronunciation{sank} \\
 & six & six & \pronunciation{see} \\
 & sept & seven & \pronunciation{set} \\
 & huit & eight & \pronunciation{weet} \\
 & neuf & nine & \pronunciation{nuhf} \\
 & dix & ten & \pronunciation{dee} \\
 & onze & eleven & \pronunciation{onz} \\
 & douze & twelve & \pronunciation{dooz} \\
 & treize & thirteen & \pronunciation{trayz} \\
 & quatorze & fourteen & \pronunciation{ka-tohrz} \\
 & quinze & fifteen & \pronunciation{kahnz} \\
 & seize & sixteen & \pronunciation{sez} \\
 & dix-sept & seventeen & \pronunciation{dee-set} \\
 & dix-huit & eighteen & \pronunciation{dee-wheet} \\
 & dix-neuf & nineteen & \pronunciation{dee-nuhf} \\
 & vingt & twenty & \pronunciation{vahn} \\
}

\section{Numbers 21-100}

\grammarexplanation[Numbers 21-100]{
To form numbers from 21 to 99, combine the tens and units:

\begin{itemize}
\item 21 = vingt-et-un (20 + 1)
\item 22 = vingt-deux (20 + 2)
\item 30 = trente
\item 35 = trente-cinq (30 + 5)
\item 40 = quarante
\item 50 = cinquante
\item 60 = soixante
\item 70 = soixante-dix (60 + 10)
\item 80 = quatre-vingts (4 × 20)
\item 90 = quatre-vingt-dix (4 × 20 + 10)
\item 100 = cent
\end{itemize}

Note: The French use a base-20 system for some numbers:
\begin{itemize}
\item 70 = 60 + 10 (soixante-dix)
\item 80 = 4 × 20 (quatre-vingts)
\item 90 = 4 × 20 + 10 (quatre-vingt-dix)
\end{itemize}
}

\section{The French Alphabet}

\vocabularylist{
02_numbers_alphabet & A & a & \pronunciation{ah} \\
 & B & b & \pronunciation{bay} \\
 & C & c & \pronunciation{say} \\
 & D & d & \pronunciation{day} \\
 & E & e & \pronunciation{uh} \\
 & F & f & \pronunciation{eff} \\
 & G & g & \pronunciation{zhay} \\
 & H & h & \pronunciation{ash} \\
 & I & i & \pronunciation{ee} \\
 & J & j & \pronunciation{zhee} \\
 & K & k & \pronunciation{ka} \\
 & L & l & \pronunciation{ell} \\
 & M & m & \pronunciation{emm} \\
 & N & n & \pronunciation{enn} \\
 & O & o & \pronunciation{oh} \\
 & P & p & \pronunciation{pay} \\
 & Q & q & \pronunciation{koo} \\
 & R & r & \pronunciation{err} \\
 & S & s & \pronunciation{ess} \\
 & T & t & \pronunciation{tay} \\
 & U & u & \pronunciation{oo} \\
 & V & v & \pronunciation{vay} \\
 & W & w & \pronunciation{doo-bluh vay} \\
 & X & x & \pronunciation{eeks} \\
 & Y & y & \pronunciation{ee-grek} \\
 & Z & z & \pronunciation{zed} \\
}

\section{Grammar: Questions with "Combien?"}

\grammarexplanation[Asking "How many?"]{
To ask "How many?" in French, use \frenchword{Combien?}:

\begin{itemize}
\item \frenchword{Combien de pommes?} = How many apples?
\item \frenchword{Combien ça coûte?} = How much does it cost?
\item \frenchword{Combien de temps?} = How much time?
\end{itemize}

Note that \frenchword{Combien} is often followed by \frenchword{de} when asking about quantity:
\begin{itemize}
\item \frenchword{Combien de livres?} = How many books?
\item \frenchword{Combien d'oranges?} = How many oranges? (using \frenchword{d'} instead of \frenchword{de} before a vowel)
\end{itemize}
}

\section{Exercises}

\exercise[Numbers Practice]{
Write the following numbers in French:

\begin{table}[h]
\centering
\begin{tabular}{l l l l l}
\toprule
15 & 27 & 43 & 68 & 99 \\
\midrule
 & & & & \\
 & & & & \\
\bottomrule
\end{tabular}
\end{table}
}

\answerkey{
15 = quinze, 27 = vingt-sept, 43 = quarante-trois, 68 = soixante-huit, 99 = quatre-vingt-dix-neuf
}

\exercise[Alphabet Practice]{
Spell the following words using the French alphabet:

\begin{itemize}
\item Your first name
\item "French"
\item "Hello"
\end{itemize}
}

\exercise[Listening Practice]{
\listening\ Listen to the numbers 0-20 being read. Write them down as you hear them. Use a text-to-speech tool with the following:

\begin{itemize}
\item zéro, un, deux, trois, quatre, cinq, six, sept, huit, neuf, dix, onze, douze, treize, quatorze, quinze, seize, dix-sept, dix-huit, dix-neuf, vingt
\end{itemize}
}

\exercise[Speaking Practice]{
\speaking\ Practice counting from 1 to 20 aloud. Then practice spelling your name.
}

\exercise[Writing Practice]{
\writing\ Write a short dialogue (6-8 lines) at a market. Use numbers and the vocabulary you've learned.
}

\section{Cultural Note}

\culturalnote{
In France, markets (\frenchword{marchés}) are an important part of daily life. You can find fresh produce, cheese, bread, and other local products. When shopping at a market:

\begin{itemize}
\item Always greet the vendor with \frenchword{Bonjour} or \frenchword{Bonsoir}
\item It's polite to ask \frenchword{Je peux vous aider?} (Can you help me?) or simply state what you want
\item Numbers are essential for specifying quantities and prices
\item Don't be afraid to ask for what you want—vendors are usually happy to help
\end{itemize}

Also, note that in French, the decimal separator is a comma, not a period. For example, 3.5 euros is written as 3,50 euros.
}

\section{Revision}

\revision{
\begin{itemize}
\item Practice counting from 0 to 100
\item Memorize the French alphabet
\item Review the numbers and alphabet vocabulary
\item Practice spelling words
\item Try the market dialogue with different items and quantities
\end{itemize}
}

\clearpage

% A1 Unit 03: Colors and Shapes
\unitheader{3}{Colors and Shapes}{Describing objects in French}

\objectives{
\item Learn to name and describe colors
\item Identify and describe shapes
\item Describe objects using colors and shapes
\item Use adjectives in French
}

\section{Dialogue: Shopping for Clothes}

\conversation{
\item \textbf{Vendeuse:} Bonjour! Je peux vous aider? \par \textit{Hello! Can I help you?}
\item \textbf{Client:} Bonjour! Je cherche une robe. \par \textit{Hello! I'm looking for a dress.}
\item \textbf{Vendeuse:} Quelle couleur? \par \textit{What color?}
\item \textbf{Client:} Bleue ou verte, s'il vous plaît. \par \textit{Blue or green, please.}
\item \textbf{Vendeuse:} Voici une robe bleue. Et celle-ci est verte. \par \textit{Here is a blue dress. And this one is green.}
\item \textbf{Client:} La bleue est très joli! Je la prends. \par \textit{The blue one is very pretty! I'll take it.}
\item \textbf{Vendeuse:} Très bien! Quelle taille? \par \textit{Very well! What size?}
\item \textbf{Client:} Taille 38. \par \textit{Size 38.}
\item \textbf{Vendeuse:} Voici. Ça fait 45 euros. \par \textit{Here you go. That's 45 euros.}
\item \textbf{Client:} Voici 50 euros. \par \textit{Here is 50 euros.}
\item \textbf{Vendeuse:} Merci! Voici votre monnaie. Au revoir! \par \textit{Thank you! Here is your change. Goodbye!}
\item \textbf{Client:} Au revoir! \par \textit{Goodbye!}
}

\section{Vocabulary}

\vocabularylist{
03_colors_shapes & bleu(e) & blue & \pronunciation{bluh} \\
 & vert(e) & green & \pronunciation{vehr} \\
 & rouge & red & \pronunciation{roozh} \\
 & jaune & yellow & \pronunciation{zhon} \\
 & blanc(he) & white & \pronunciation{blon/blonch} \\
 & noir(e) & black & \pronunciation{nwahr} \\
 & gris(e) & gray & \pronunciation{gree} \\
 & rose & pink & \pronunciation{roz} \\
 & orange & orange & \pronunciation{oh-rahnzh} \\
 & violet(te) & purple & \pronunciation{vee-oh-lay} \\
 & marron & brown & \pronunciation{ma-ron} \\
 & beige & beige & \pronunciation{bezh} \\
}

\vocabularylist{
03_colors_shapes & un cercle & a circle & \pronunciation{uhn sehr-kluh} \\
 & un carré & a square & \pronunciation{uhn ka-ray} \\
 & un triangle & a triangle & \pronunciation{uhn tree-ong-gluh} \\
 & un rectangle & a rectangle & \pronunciation{uhn rec-tong-gluh} \\
 & un ovale & an oval & \pronunciation{uhn oh-val} \\
 & une ligne & a line & \pronunciation{oon leen-yuh} \\
 & rond(e) & round & \pronunciation{ron} \\
 & pointu(e) & pointed & \pronunciation{pwan-too} \\
 & grand(e) & big & \pronunciation{grahn} \\
 & petit(e) & small & \pronunciation{puh-tee} \\
}

\section{Grammar: Adjectives}

\grammarexplanation[Adjective Agreement]{
In French, adjectives must agree in gender and number with the nouns they describe:

\begin{itemize}
\item \frenchword{Un livre bleu} = A blue book (masculine singular)
\item \frenchword{Une robe bleue} = A blue dress (feminine singular)
\item \frenchword{Des livres bleus} = Blue books (masculine plural)
\item \frenchword{Des robes bleues} = Blue dresses (feminine plural)
\end{itemize}

\textbf{Rules for adjective agreement:}
\begin{itemize}
\item Most adjectives add \frenchword{-e} for feminine: \frenchword{bleu} → \frenchword{bleue}
\item Most adjectives add \frenchword{-s} for plural: \frenchword{bleu} → \frenchword{bleus}
\item For feminine plural, add \frenchword{-es}: \frenchword{bleu} → \frenchword{bleues}
\item Some adjectives have irregular feminine forms
\end{itemize}

\textbf{Position of adjectives:}
Most adjectives come AFTER the noun in French:
\begin{itemize}
\item \frenchword{Une voiture rouge} = A red car
\item \frenchword{Un ciel bleu} = A blue sky
\end{itemize}

However, some common adjectives come BEFORE the noun:
\begin{itemize}
\item \frenchword{Un grand arbre} = A tall tree
\item \frenchword{Une petite maison} = A small house
\item \frenchword{Un bon livre} = A good book
\item \frenchword{Une belle journée} = A beautiful day
\end{itemize}
}

\section{Grammar: Describing Objects}

\grammarexplanation[Using "C'est" and "Ce sont"]{
To describe objects, use \frenchword{c'est} (it is) or \frenchword{ce sont} (they are):

\begin{itemize}
\item \frenchword{C'est une table.} = It is a table.
\item \frenchword{C'est un livre intéressant.} = It is an interesting book.
\item \frenchword{Ce sont des chaises.} = They are chairs.
\item \frenchword{Ce sont des fleurs belles.} = They are beautiful flowers.
\end{itemize}

To describe colors and shapes:
\begin{itemize}
\item \frenchword{C'est bleu.} = It is blue.
\item \frenchword{C'est rond.} = It is round.
\item \frenchword{Ce sont verts.} = They are green.
\item \frenchword{Ce sont carrés.} = They are square.
\end{itemize}
}

\section{Exercises}

\exercise[Colors Practice]{
Match the French color words with their English meanings:

\begin{table}[h]
\centering
\begin{tabular}{l l}
\toprule
\textbf{French} & \textbf{English} \\
\midrule
1. bleu & a. green \\
2. vert & b. red \\
3. rouge & c. blue \\
4. jaune & d. yellow \\
5. noir & e. black \\
\bottomrule
\end{tabular}
\end{table}
}

\answerkey{
1. c, 2. a, 3. b, 4. d, 5. e
}

\exercise[Shapes Practice]{
Draw the following shapes and label them in French:

\begin{itemize}
\item un cercle
\item un carré
\item un triangle
\item un rectangle
\end{itemize}
}

\exercise[Description Practice]{
Describe the following objects using colors and shapes:

\begin{itemize}
\item A red circle
\item A blue square
\item A green triangle
\item A yellow rectangle
\end{itemize}
}

\answerkey{
Un cercle rouge, Un carré bleu, Un triangle vert, Un rectangle jaune
}

\exercise[Listening Practice]{
\listening\ Listen to the following colors and shapes. Write them down in French:

\begin{itemize}
\item blue, green, red, yellow, black
\item circle, square, triangle, rectangle
\end{itemize}
}

\exercise[Speaking Practice]{
\speaking\ Describe your favorite color and a shape that represents it. Use complete sentences.
}

\exercise[Writing Practice]{
\writing\ Write a short paragraph (4-6 sentences) describing your favorite outfit. Include colors and any patterns or shapes.
}

\section{Cultural Note}

\culturalnote{
In France, fashion and style are important aspects of culture. The French are known for their sense of style and quality clothing. When describing clothes or other items:

\begin{itemize}
\item Colors are often used to describe moods and impressions
\item The French appreciate attention to detail in design
\item Quality is often valued over quantity
\end{itemize}

Also, note that in French, the word for "fashion" is \frenchword{la mode}, and Paris is considered one of the fashion capitals of the world.
}

\section{Revision}

\revision{
\begin{itemize}
\item Practice naming colors and shapes
\item Review adjective agreement rules
\item Practice describing objects using colors and shapes
\item Try the clothing dialogue with different items and colors
\end{itemize}
}

\clearpage

% Continue with remaining A1 units...
% This is a sample of the first 3 units of A1
% The complete book will have all 6 levels with multiple units each

\chapter{Remaining A1 Units}
\label{chap:a1-remaining}

\section*{A1 Unit 04: Family and Friends}
\section*{A1 Unit 05: House and Objects}
\section*{A1 Unit 06: Food and Drinks}
\section*{A1 Unit 07: Clothing and Fashion}
\section*{A1 Unit 08: Hobbies and Free Time}
\section*{A1 Unit 09: Daily Routines}

\begin{center}
\includesvg[height=5cm]{mascot_amelie_encouraging}
\end{center}

\begin{notebox}
The remaining A1 units will follow the same comprehensive structure with dialogues, vocabulary, grammar explanations, exercises, cultural notes, and revision sections.
\end{notebox}

% ============================================================================
% LEVEL A2 - ELEMENTARY
% ============================================================================

\part{A2: Elementary}
\label{part:a2}

\chapter{A2: Building on the Basics}
\label{chap:a2}

\section*{About Level A2}
\markboth{Level A2: Building on the Basics}{Level A2: Building on the Basics}

\levelAtwo\ At the end of Level A2, you will be able to:

\begin{itemize}
\item Understand sentences and frequently used expressions related to areas of most immediate relevance
\item Communicate in simple and routine tasks requiring a simple and direct exchange of information
\item Describe in simple terms aspects of your background, immediate environment, and matters in areas of immediate need
\end{itemize}

\begin{center}
\includesvg[height=5cm]{mascot_amelie_explaining}
\end{center}

\section*{A2 Units}
\begin{itemize}
\item Unit 01: Travel and Transport
\item Unit 02: Hotel and Restaurant
\item Unit 03: Shopping
\item Unit 04: Directions and Places
\item Unit 05: Weather and Seasons
\item Unit 06: Time and Dates
\item Unit 07: Work and Professions
\item Unit 08: Health and Body
\item Unit 09: Education and Learning
\item Unit 10: Technology and Communication
\end{itemize}

% ============================================================================
% LEVEL B1 - INTERMEDIATE
% ============================================================================

\part{B1: Intermediate}
\label{part:b1}

\chapter{B1: Developing Independence}
\label{chap:b1}

\section*{About Level B1}
\markboth{Level B1: Developing Independence}{Level B1: Developing Independence}

\levelBone\ At the end of Level B1, you will be able to:

\begin{itemize}
\item Understand the main points of clear standard input on familiar matters regularly encountered in work, school, leisure, etc.
\item Deal with most situations likely to arise while traveling in an area where the language is spoken
\item Produce simple connected text on topics that are familiar or of personal interest
\item Describe experiences and events, dreams, hopes, and ambitions and briefly give reasons and explanations for opinions and plans
\end{itemize}

\section*{B1 Units}
\begin{itemize}
\item Unit 01: Past Experiences
\item Unit 02: Future Plans
\item Unit 03: Opinions and Arguments
\item Unit 04: Media and News
\item Unit 05: Arts and Culture
\item Unit 06: Society and Traditions
\item Unit 07: Environment and Nature
\item Unit 08: Science and Technology
\item Unit 09: Global Issues
\item Unit 10: Personal Development
\end{itemize}

% ============================================================================
% LEVEL B2 - UPPER INTERMEDIATE
% ============================================================================

\part{B2: Upper Intermediate}
\label{part:b2}

\chapter{B2: Achieving Fluency}
\label{chap:b2}

\section*{About Level B2}
\markboth{Level B2: Achieving Fluency}{Level B2: Achieving Fluency}

\levelBtwo\ At the end of Level B2, you will be able to:

\begin{itemize}
\item Understand the main ideas of complex text on both concrete and abstract topics, including technical discussions in your field of specialization
\item Interact with a degree of fluency and spontaneity that makes regular interaction with native speakers quite possible without strain for either party
\item Produce clear, detailed text on a wide range of subjects and explain a viewpoint on a topical issue giving the advantages and disadvantages of various possibilities
\end{itemize}

\section*{B2 Units}
\begin{itemize}
\item Unit 01: Advanced Conversation
\item Unit 02: Professional Communication
\item Unit 03: Literary Analysis
\item Unit 04: Current Affairs
\item Unit 05: Philosophical Discussions
\item Unit 06: Business French
\item Unit 07: Academic Writing
\item Unit 08: Cultural Critique
\item Unit 09: Debates and Discussions
\item Unit 10: Advanced Grammar
\end{itemize}

% ============================================================================
% LEVEL C1 - ADVANCED
% ============================================================================

\part{C1: Advanced}
\label{part:c1}

\chapter{C1: Mastering Complexity}
\label{chap:c1}

\section*{About Level C1}
\markboth{Level C1: Mastering Complexity}{Level C1: Mastering Complexity}

\levelCone\ At the end of Level C1, you will be able to:

\begin{itemize}
\item Understand a wide range of demanding, longer texts, and recognize implicit meaning
\item Express yourself fluently and spontaneously without much obvious searching for expressions
\item Use language flexibly and effectively for social, academic, and professional purposes
\item Produce clear, well-structured, detailed text on complex subjects, showing controlled use of organizational patterns, connectors, and cohesive devices
\end{itemize}

\section*{C1 Units}
\begin{itemize}
\item Unit 01: Advanced Reading
\item Unit 02: Nuanced Writing
\item Unit 03: Complex Grammar
\item Unit 04: Specialized Vocabulary
\item Unit 05: Literary French
\item Unit 06: Professional Presentations
\item Unit 07: Cultural Analysis
\item Unit 08: Advanced Translation
\end{itemize}

% ============================================================================
% LEVEL C2 - PROFICIENCY
% ============================================================================

\part{C2: Proficiency}
\label{part:c2}

\chapter{C2: Near-Native Proficiency}
\label{chap:c2}

\section*{About Level C2}
\markboth{Level C2: Near-Native Proficiency}{Level C2: Near-Native Proficiency}

\levelCtwo\ At the end of Level C2, you will be able to:

\begin{itemize}
\item Understand with ease virtually everything heard or read
\item Summarize information from different spoken and written sources, reconstructing arguments and accounts in a coherent presentation
\item Express yourself spontaneously, very fluently, and precisely, differentiating finer shades of meaning even in more complex situations
\end{itemize}

\section*{C2 Units}
\begin{itemize}
\item Unit 01: Mastery of All Skills
\item Unit 02: Specialized Fields
\item Unit 03: Advanced Literature
\item Unit 04: Professional Mastery
\item Unit 05: Cultural Expertise
\item Unit 06: Advanced Translation
\item Unit 07: Creative Writing
\item Unit 08: Comprehensive Review
\end{itemize}

% ============================================================================
% REFERENCE SECTIONS
% ============================================================================

\part{Reference Material}
\label{part:reference}

\chapter{Grammar Index}
\label{chap:grammar}

\section*{A1 Level Grammar}
\section*{A2 Level Grammar}
\section*{B1 Level Grammar}
\section*{B2 Level Grammar}
\section*{C1 Level Grammar}
\section*{C2 Level Grammar}

\chapter{Glossary}
\label{chap:glossary}

French-English and English-French glossary with all vocabulary from the book.

\chapter{Verb Conjugation Tables}
\label{chap:verbs}

Complete conjugation tables for all verbs introduced in the book.

\chapter{Pronunciation Guide}
\label{chap:pronunciation}

Comprehensive guide to French pronunciation with examples and practice exercises.

% ============================================================================
% ANSWER KEYS
% ============================================================================

\part{Answer Keys}
\label{part:answers}

\chapter{A1 Level Answer Keys}
\label{chap:answers-a1}
\chapter{A2 Level Answer Keys}
\label{chap:answers-a2}
\chapter{B1 Level Answer Keys}
\label{chap:answers-b1}
\chapter{B2 Level Answer Keys}
\label{chap:answers-b2}
\chapter{C1 Level Answer Keys}
\label{chap:answers-c1}
\chapter{C2 Level Answer Keys}
\label{chap:answers-c2}

% ============================================================================
% STUDY GUIDE
% ============================================================================

\part{Study Guide}
\label{part:studyguide}

\chapter{Learning Strategies}
\label{chap:strategies}
\chapter{Assessment Guides}
\label{chap:assessment}
\chapter{Resource Recommendations}
\label{chap:resources}

\end{document}
